\section{lista 4 draft só soluções}

\begin{exercise}[1.b]
\end{exercise}

\begin{proof}

	Para a primeira equação, seja $Z = \Ex[X \mid X \land a]$, vamos
	mostrar que se $\omega \in \{X \leq a\}$, então $Z(\omega) = X$ q.c e
	se $\omega \in \{X > a\}$, então $Z(\omega) > a$ q.c. Seja
	\[B = \{X \leq a\} \cap \{X > Z\} = \{X \land a \leq a\} \cap \{X
	\land a > Z\} \in \sigma(X \land a).\]
	Pela propriedade da esperança condicional e pela definição de $B$, segue que
	\[
		0 \leq \Ex[\id[B](X - Z)] = \Ex[\id[B](X \land a)] - \Ex[\id[B](Z)] = 0
	\]
	portanto $\Pr(B) = 0$. Da mesma forma, considerando $B' = \{X \leq
	a\} \cap \{X < Z\}$, chegamos a conclusão que $\Pr(B') = 0$. Segue
	que $X = Z$ q.c em $\{X \leq a\}$.

	Agora seja $C \subset \{X \land a = a\}$ qualquer conjunto $\sigma(X
	\land a)$ mensurável. Notamos que
	\[
		\Ex[\id[C](Z - a)] = \Ex[\id[C]X - \id[C]a] \geq \Ex[\id[C](a - a)] = 0
	\]
	logo $Z - a \geq 0$ q.c em $\{X \geq a\}$.

	A primeira equação segue dessa separação em casos. Se $X \leq a$,
	então $Z = X$ q.c e portanto $Z \land a = X \land a$. Se $X \geq a$,
	então $Z \geq a$ q.c
	e temos $Z \land a = a = X \land a$ q.c.

	A segunda e terceira equações são simétricas em $X$ e $Y$, então
	basta mostrar a segunda. Vamos mostrar que $Y \land a$ satisfaz a
	condição da experança condicional.
	Seja $C \in \sigma(Y \land a)$ e vamos dividir $C = C^- \sqcup C^+$ com
	\begin{align*}
		C^- = (C \cap \{Y < a\}) = (C \cap \{Y \land a < a\}) \in \sigma(Y \land a)\\
		C^+ = (C \cap \{Y \geq a\}) = (C \cap \{Y \land a = a\}) \in \sigma(Y \land a)
	\end{align*}
	Temos
	\begin{equation}
		\label{lista4:ex1:eq1}
		\Ex[\id[C] Y \land a] = \Ex[\id[C^-]Y] + \Ex[\id[C^+] a].
	\end{equation}
	Pela hipótese $\Ex[Y \mid X] = X$ e sabendo que $\sigma(X) = \sigma(Y)$, vemos
	\[
		\Ex[\id[C^-]Y] = \Ex[ \id[C^-]\Ex[Y \mid X]] = \Ex[\id[C^-]X] =
		\Ex[\id[C^-](X \land a)].
	\]
	Falta provar que $\Ex[\id[C^+]a] = \Ex[\id[C^+](X \land a)]$.
	Veremos que em $\{Y \geq a\}$ vale que $X \geq a$ q.c e a equação
	segue imediatamente.
	Considere $D = \{Y \geq a\} \cap \{X < a\}$ e
	\[
		0 \leq \Ex[\id[D]Y - X] = \Ex[\id[D]Y] - \Ex[\id[D]X] = 0.
	\]
	Segue que $\Pr(D) = 0$ e portanto
	$\Ex[\id[C^+]X \land a] = \Ex[\id[C^+]a]$. Juntando esses dois
	resultados na equação [\ref{lista4:ex1:eq1}], temos
	\[
		\Ex[\id[C] Y \land a] = \Ex[\id[C] X \land a]
	\]
	como queríamos mostrar.

	Falta finalizar a questão. Vamos mostrar que $\Pr(X < Y) = 0$, por
	simetria $\Pr(X = Y) = 1$. Escrevemos crucialmente que
	\[
		U := \{X < Y\} = \bigcup_{a \in \mathbb{Q}} \{X \land a < Y \land a\}
	\]
	Em particular, dado $s \in \mathbb{Q}$, $U_s = \bigcup_{a \in
	\mathbb{Q}, a \leq s}\{X \land a < Y \land a\} \in \sigma(X \land a)
	= \sigma(Y \land a)$. E temos que
	$U_s$ cresce para $U$.

	Vamos ver que $\Ex[\id[U](Y - X)] = 0$. Seja $\eps > 0$, existe
	$\delta > 0$ tal que se $\Pr(A) < \delta$, então $\Ex[\id[A] (|X| +
	|Y|)] < \eps$ (pois $X,Y \in L^1$).
	Tome $s$ tão grande que $\Pr(U\setminus U_s) < \delta$. Segue que,
	aplicando o TCD,
	\begin{align*}
		\Ex[\id[U](Y - X)] \leq \Ex[\id[U_s] (Y - X)] + \eps = \Ex[\id[U_s]
		\lim_{a \to \infty} (Y\land a - X\land a)] + \eps\\
		= \lim_{a \to \infty} \Ex[\id[U_s] Y\land a] - \Ex[\id[U_s] X \land
		a] + \eps = \eps.
	\end{align*}
	Como $\eps$ foi arbitrário, $\Ex[\id[U](Y - X)] \leq 0$ e segue o resultado.
\end{proof}

\begin{exercise}[5]
\end{exercise}

\begin{proof}
	(a) Para encontrar a distribuição de $X$, basta calcular
	\[
		\Pr(X = n) = \lim_{t \to \infty} \Pr(X = n, Y \leq t) = b
		\int_{0}^\infty \frac{(ay)^n}{n!}e^{-(a+b)y}dy.
	\]
	Lembrando da função $\Gamma$ e do fatorial e fazendo a substituição
	$(a+b)y = u$, é fácil ver que
	\[
		b \int_{0}^\infty \frac{(ay)^n}{n!}e^{-(a+b)y} dy =
		\frac{ba^n}{n!}\frac{1}{(a + b)^{n+1}} \int_0^\infty u^n e^{-u} du
		= \bigg(\frac{a}{a + b}\bigg)^n \frac{b}{(a + b)}
	\]
	que é uma geométrica que conta quantidade de fracassos com parâmetro
	de sucesso $p = b/(a+b)$.

	Para determinar $Y$, basta usar Tonelli-Fubini e vemos que
	\[
		\Pr(Y \leq t) = \sum_{n \geq 0} b\int_0^t
		\frac{(ay)^n}{n!}e^{-(a+b)y} = \int_0^t be^{-(a+b)y} \sum_{n \geq
		0} \frac{(ay)^n}{n!} dy = \int_0^\infty be^{-by}dy
	\]
	que é uma exponêncial com parâmetro $b$.

	(b) Segue da completamente da expressão que (avaliando em cada ponto $n$)
	\[
		\Pr(Y \leq t \mid X) = b\int_0^t \frac{(ay)^X}{X!}e^{-(a+b)y} dy.
	\]

	(c) Fixado $X$, temos (usando a mesma substituição que antes e $\Gamma$)
	\[
		\Ex [Y \mid X] = b\int_0^\infty \frac{a^Xy^{X + 1}}{X!}e^{-(a+b)y}dy =
		\frac{(X+1)!}{X!} \bigg(\frac{a}{a + b}\bigg)^X \frac{b}{(a + b)^2}
		= \frac{(X + 1)a^Xb}{(a+b)^{X + 2}}.
	\]
	Logo, condicionando primeiro em $X$,
	\[
		\Ex[Y/(X+1)] = \Ex[ \Ex[Y/(X+1) \mid X]] = \Ex[ \Ex[Y \mid
		X]/(X+1)] = \frac{b}{(a+b)^2}\Ex[a^X/(a + b)^{X}]
	\]
	usando a distribuição de $X$ temos
	\[
		\Ex[a^X/(a + b)^{X}] = \sum_{n\geq 0} \Pr(X = n)
		\frac{a^n}{(a+b)^n} = \frac{b}{(a + b)} \sum_{n \geq 0}
		\bigg(\frac{a}{a + b}\bigg)^{2n}
	\]
	somando a geométrica achamos
	\[
		\Ex[a^X/(a + b)^{X}] = \frac{b}{(a + b)} \frac{1}{1 - (a/(a+b))^2}
		= \frac{a+b}{2a + b}.
	\]
	Por fim
	\[
		\Ex[Y/(X+1)] = \frac{b}{(a+b)(2a + b)}.
	\]

	(d) Vamos tentar o palpite óbvio do que deveria ser supondo. Se o
	mundo fosse finito e bonito, teríamos
	\[
		\Ex[\id[X = n] \mid Y = y] \approx \Pr(X = n \mid Y = y) \approx
		\frac{\Pr(X = n \land Y = y)}{\Pr(Y = y)}.
	\]
	Como não é o caso, vamos nos guiar por isso para encontrar um
	candidato para $\Ex[\id[X = n] \mid Y]$.
	É intuitivo ver que $\Pr(Y = y) \sim dY(y) = be^{-by}$ pela letra
	(a). Seria natural que
	\[
		\frac{\Pr(X = n \land Y = y)}{\Pr(Y = y)} \sim
		\frac{b(ay)^ne^{-(a+b)y}}{n!} \frac{1}{be^{-by}} = \frac{b(ay)^n e^{-ay}}{n!}
	\]
	Seja então
	\[
		Z_n = \frac{b(aY)^n e^{-aY}}{n!}.
	\]
	Vamos ver que $Z_n = \Ex[\id[X = n] \mid Y]$ pois satisfaz a
	propriedade da condicional. Como $Y$ é uma exponêncial
	e $e^{aY} \leq 1$ vale que $||Z_n||_{L^1} \leq ba^n||Y||_{L^n} <
	\infty$. Logo $Z_n \in L^1$. Vamos provar que
	\[\Ex[\id[C] \id[X = n]] = \Ex[\id[C]Z_n]\]
	quando $C$ é da forma $\{s < Y < t\}$ que é um $\Pi$-sistema gerador
	de $\sigma(Y)$ o que conclui a afirmação.
	Mas isso segue da seguinte conta
	\begin{align*}
		\Ex[\id[C] \id[X = n]] = \Pr(X = n, s < Y < t) = b\int_s^t
		\frac{(ay)^n}{n!}e^{-(a+b)y}dy\\
		= b\int_s^t \frac{b(ay)^n e^{-ay}}{n!} dY(y) = \Ex[\id[C] Z_n].
	\end{align*}

	Para (e), basta então considerar $X = \sum_{n \geq 1} n\id[X = n]$,
	donde por linearidade e convergência monótona segue que
	\[
		\Ex[X | Y] = \sum_{n \geq 1} n Z_n = b e^{-aY}\sum_{n \geq 1}
		\frac{(aY)^n}{(n-1)!} = baY.
	\]
\end{proof}

\begin{exercise} (9)
\end{exercise}

\begin{proof}
	Vamos mudar o paradigma. Enumeramos os passageiros por $[n]$ e
	"seus" assentos por $[n]$ também. Esqueçamos que o assento $1$ foi designado ao
	passageiro $1$, vamos somente guardar a informação que os outros
	passageiros $\{2,\dots, n\}$ nunca tentarão primeiramente sentar no
	assento $1$.
	Essa forma de pensar ajuda a justificar a indução.

	Para $1 \leq i \leq n$, seja $X_i$ o assento em que o $i$-ésimo
	passageiro sentou.  Vamos provar por indução que se $n \geq 2$,
	$\Pr(X_n = n) = 1/2$.
	É óbvio que se $n = 1$, então $\Pr(X_1 = 1) = 1$.

	Caso $n = 2$. Esse é bem simples, ou $X_1 = 1$ e sobra somente $X_2
	= 2$ ou no outro caso $X_1 = 2$ e $X_2 = 1$. Como cada caso tem
	probabilidade $1/2$, segue que
	$\Pr(X_2 = 2) = 1/2$.

	Caso $n > 2$. Vamos escrever $\Pr(X_n = n)$ usando a condicional,
	\[
		\Pr(X_n = n) = \sum_{j = 1}^n \Pr(X_n = n \mid X_1 = j) \Pr(X_1 =
		j) = \sum_{j = 1}^n \frac{\Pr(X_n = n \mid X_1 = j)}{n}.
	\]
	Agora vamos analisar o que é $\Pr(X_n = n \mid X_1 = j)$. Claramente,
	se $X_1 = 1$, todos os próximos passageiros sentarão em seus
	respectivos lugares sem achar ninguém
	os ocupando e o processo segue deterministicamente, portanto segue
	que $\Pr(X_n = n \mid X_1 = 1) = 1$. Se $X_1 = n$, como estará para
	sempre ocupado, é impossível que
	$X_n = n$ e logo temos $\Pr(X_n = n \mid X_1 = n) = 0$.

	O interessante é o que acontece no meio. Se $X_1 = j$ com $1 < j <
	n$, então temos
	que o primeiro passageiro sentou-se no assento $j$ e todos os
	passageiros $2, 3 , \dots, j-1$ sentam-se normalmente em seus
	respectivos assentos.
	Quando o passageiro $j$ embarca, ele vê que alguém se sentou no
	lugar $j$ e em vez de sentar em um lugar aleatório por raiva, ele
	volta ao início do avião e acredita que lembrou
	errado do seu número e esqueceu qual era. Agora o $j$-ésimo
	passageiro se encontra na mesma situação que o primeiro estava, ele
	não sabe seu assento, vai sentar em um aleatório
	e tem um assento específico (o $1$) que nenhum dos passageiros $j+1,
	j+2 \dots n$ estão interessados. Como esse é exatamente o mesmo
	modelo que estamos lidando (mas com menos passageiros),
	por indução segue que $\Pr(X_n = n | X_1 = j) = 1/2$. Substituindo, temos
	\[
		\Pr(X_n = n) = \frac{1}{n} + \sum_{j = 1}^{n-1} \frac{1}{2n} = \frac{1}{2}.
	\]

	% Vamos provar por indução e mudaremos o paradigma também.
	% Esqueçamos que o primeiro passageiro tenha um assento designado,
	% basta reinterpretar que nenhum dos restantes passagereiros tentará
	% o primeiro assento. caso ele o escolha, todos os outros
	% passageiros sentarão normalmente.
\end{proof}
